% Nginx en contenedor (router-nginx)
\subsection{Nginx en contenedor (router-nginx)}\label{sec:7_1_router_nginx}

\subsubsection{Contenedor router-nginx con -{}-network host}
\texttt{router-nginx} es un nginx oficial (\texttt{nginx:alpine}) que funciona como \emph{proxy inverso}: recibe las peticiones de los usuarios, termina la conexión HTTPS y reenvía cada petición al frontend o al microservicio que corresponde según la ruta (sección \ref{sec:7_4_defensoria_conf}). No tiene imagen propia: su comportamiento lo define por completo el archivo \texttt{defensoria.conf}, que se monta desde la VPS.

A diferencia de los demás contenedores, se crea con \texttt{-{}-network host}: no tiene red propia, sino que usa directamente la red de la VPS. Esto permite que:
\begin{itemize}
  \item Escuche en los puertos 80 y 443 de la VPS sin necesidad de publicarlos con \texttt{-p}.
  \item Alcance los cinco frontends por \texttt{127.0.0.1:8090}--\texttt{8095}, que son los puertos que estos publican en la VPS.
  \item Registre la IP real de cada visitante en sus logs, en lugar de una IP interna de Podman.
\end{itemize}
El contenedor se crea en la sección \ref{sec:7_3_certificados_certbot}, una vez obtenido el certificado.

\subsubsection{Estructura de /apps/aplicaciones/defensoria/router (config y certbot-webroot)}
\begin{lstlisting}
/apps/aplicaciones/defensoria/router/
|-- config/
|   `-- defensoria.conf      -> /etc/nginx/conf.d/default.conf (solo lectura)
`-- certbot-webroot/         -> /var/www/certbot (solo lectura)
/etc/letsencrypt/            -> /etc/letsencrypt (solo lectura, certificados)
\end{lstlisting}
\begin{itemize}
  \item \texttt{config/defensoria.conf}: configuración de nginx, copiada en la sección \ref{sec:4_4_transferencia_vps}.
  \item \texttt{certbot-webroot/}: carpeta donde Certbot deja el archivo temporal del reto de Let's Encrypt y desde donde nginx lo sirve.
  \item \texttt{/etc/letsencrypt}: carpeta del sistema donde Certbot guarda los certificados; se monta en el contenedor para que nginx los lea.
\end{itemize}
Los tres volúmenes se montan en solo lectura (\texttt{:ro}): nginx solo necesita leerlos.
\begin{lstlisting}
# find /apps/aplicaciones/defensoria/router -printf "%M %u %g %p\n"
\end{lstlisting}
\captura{front-01-router-estructura}{Carpeta router con defensoria.conf y certbot-webroot}
